\color{unimain}{\subsection{Направленная высокочастотная защита (НВЧЗ)}}
\color{black}

\begin{enumerate}

%%%%%%%%%%%%%%%%%%%%%%%%%%%%%%%%%%%%%%%%%%%%%%%%%% ОБШИЕ СВЕДЕНИЯ %%%%%%%%%%%%%%%%%%%%%%%%%%%%%%%%%%%%%%%%%%%%%%%%%%
\item Общие сведения

\begin{enumerate}

\item
Направленная высокочастотная защита (НВЧЗ) является быстродействующей защитой линии с абсолютной селективностью и выполняет функцию основной защиты линии от всех видов КЗ. 

\item
НВЧЗ реализуется в виде полукомплектов, устанавливаемых на концах линии, каждый из которых включает микропроцессорный терминал релейной защиты и комплект оборудования высокочастотной (ВЧ) связи, обеспечивающий формирование и приём ВЧ-сигналов по каналу связи с использованием внешних приёмопередатчиков (ПП).

\item 
НВЧЗ селективно срабатывает:
\begin{itemize}
    \item при всех видах КЗ, в том числе при постановке линии под напряжение; 
    \item при переходе внешнего КЗ во внутреннее, в том числе при реверсе мощности.
\end{itemize}

\item 
Обеспечивается несрабатывание защиты:
\begin{itemize} 
\item при постановке линии под напряжение и включении в транзит при отсутствии КЗ;
\item при внешних КЗ, в том числе при реверсе мощности;
\item при синхронных качаниях и в режиме асинхронного хода;
\item при КЗ за трансформаторами ответвлений (отпаек) на линиях с ответвлениями;
\item при бросках тока намагничивания трансформаторов без КЗ;
\item при несимметрии токов, вызванной тяговой нагрузкой;
\item при внешних КЗ с насыщением ТТ (при соблюдении заявленных производителем требований к ТТ).
\end{itemize}

\item
Принцип действия НВЧЗ основан на сравнении направлений мощности по концам защищаемой линии, определяемых органом направления мощности обратной последовательности (ОНМОП) при несимметричных КЗ и направленным реле сопротивления при симметричных.
Сравнение направлений мощности осуществляется косвенно – посредством обмена ВЧ-сигналами между полукомплектами. 

\item
НВЧЗ содержит две основные группы измерительных органов (ИО): блокирующие – действующие на пуск ВЧ-передатчика, и отключающие – действующие на отключение выключателя.
В начальный момент повреждения срабатывают блокирующие органы, обеспечивая ускоренный запуск ВЧ-передатчика и передачу блокирующего ВЧ-сигнала на противоположный конец линии.
Место КЗ, внутри или вне защищаемой линии, определяется по отсутствию или наличию постоянного ВЧ-сигнала в канале связи. 

\item
Если повреждение находится в пределах защищаемой линии, то после срабатывания отключающих органов, имеющих более грубые уставки, ВЧ-передатчики каждого полукомплекта останавливаются, при этом блокирующий ВЧ-сигнал снимается, и оба полукомплекта действуют на отключение своих выключателей.
Если повреждение находится вне защищаемой линии, «за спиной» одного из полукомплектов, то в этом полукомплекте отключающие органы не срабатывают, его ВЧ-передатчик остается в работе, и блокирующий ВЧ-сигнал продолжает передаваться по каналу связи, предотвращая срабатывание защиты обоих полукомплектов.

\item
Команда на пуск ВЧ-передатчика (логический сигнал \textit{«Пуск ВЧ»}) формируется при срабатывании блокирующих органов, отстроенных от нагрузочного режима.
Блокирующие органы реализованы с возможностью независимой настройки их уставок по:
\begin{itemize} 
\item току обратной последовательности \textit{«ИО I2 блок НВЧЗ»};
\item приращению тока прямой последовательности \textit{«ИО dI1 блок НВЧЗ»};
\item приращению тока обратной последовательности \textit{«ИО dI2 блок НВЧЗ»};
\item напряжению обратной последовательности \textit{«ИО U2 блок НВЧЗ»};
\item сопротивлению «РС блок НВЧЗ».
\end{itemize}

\item
Подготовка цепи отключения осуществляется от отключающих органов.
Отключающие органы реализованы с возможностью независимой настройки их уставок по:
\begin{itemize} 
\item току обратной последовательности \textit{«ИО I2 откл НВЧЗ»};
\item приращению тока прямой последовательности \textit{«ИО dI1 откл НВЧЗ»};
\item приращению тока обратной последовательности \textit{«ИО dI2 откл НВЧЗ»};
\item напряжению обратной последовательности \textit{«ИО U2 откл НВЧЗ»};
\item сопротивлению \textit{«РС откл НВЧЗ»};
\item направлению мощности обратной последовательности \textit{«ОНМОП откл НВЧЗ»};
\item \textls[-30]{току обратной последовательности с торможением от тока прямой последовательности \textit{«ИО I2т откл НВЧЗ»}.}
\end{itemize}

\item
Измерительные органы НВЧЗ представлена на рисунке \ref{pict_nvchz:IO_NVChZ}.

\medskip

\begin{figure}[H]
\centering
\includegraphics[width=0.94\textwidth,height=0.94\textheight,keepaspectratio]{\fbpath/01.01 ВЧ/561.0771 НВЧЗ/img/IO_NVChZ.pdf}
\captionof{figure}{Измерительные органы НВЧЗ}
\label{pict_nvchz:IO_NVChZ}
\end{figure}

\medskip

\item 
\textls[-15]{Уставочные параметры измерительных органов НВЧЗ представлены в таблице \ref{app_set:nvchz}.}

\end{enumerate}


%%%%%%%%%%%%%%%%%%%%%%%%%%%%%%%%%%%%%%%%%%%%%%%%%% БЛОКИРУЮЩИЕ ИО %%%%%%%%%%%%%%%%%%%%%%%%%%%%%%%%%%%%%%%%%%%%%%%%%%
\item Блокирующие измерительные органы

\begin{enumerate}

\item
Блокирующие органы по приращениям токов прямой и обратной последовательностей рекомендуется применять в сетях с тяговой нагрузкой, для которых характерны значительные небалансы токов обратной последовательности.
Время возврата этих органов составляет 250 мс. 

\item
Блокирующий орган по сопротивлению (блокирующее РС) имеет обратнонаправленную характеристику срабатывания с охватом начала координат, которая задается уставками \textbf{«Rср РС блок НВЧЗ»}, \textbf{«Хср РС блок НВЧЗ»} и \textbf{«Ф1 РС блок НВЧЗ»}.
Для отстройки от нагрузки предусмотрен вырез в характеристике срабатывания, который определяется уставками \textbf{«Rнг РС блок НВЧЗ»} и \textbf{«Фнг РС блок НВЧЗ»}.
Работа блокирующего РС может быть заблокирована при выявлении неисправности в цепях напряжения при введенном программном ключе \textbf{«БНН РС блок НВЧЗ»}.

\item
Характеристика срабатывания блокирующего РС представлена на рисунке \ref{pict_nvchz:RS_blk_NVChZ}.

\medskip

\begin{figure}[H]
\centering
\includegraphics[width=0.35\textwidth,height=0.35\textheight,keepaspectratio]{\fbpath/01.01 ВЧ/561.0771 НВЧЗ/img/RS_blk_NVChZ.pdf}
\captionof{figure}{Характеристика срабатывания блокирующего РС НВЧЗ}\label{pict_nvchz:RS_blk_NVChZ}
\end{figure}

\medskip

\item
При подключении параллельных линий к общим шинам на обоих концах НВЧЗ может сработать неселективно из-за реверса мощности.
В начальный момент повреждения на параллельной линии (линия W2) срабатывают блокирующие ИО полукомплектов А1 и А2, обеспечивая ускоренный запуск ВЧ-передатчиков.
Затем срабатывают отключающие ИО полукомплекта А2, но не срабатывают отключающие ИО полукомплекта А1, его ВЧ-передатчик остается в работе – наличие ВЧ-сигнала в канале связи предотвращает срабатывание НВЧЗ (см. рисунок \ref{pict_nvchz:REV_S_a}). 

\medskip

\begin{figure}[H]
\centering
\begin{subfigure}{0.47\textwidth}
    \centering
    \includegraphics[width=\textwidth]{\fbpath/01.01 ВЧ/561.0771 НВЧЗ/img/REV_S_a.pdf}
    \caption{а) начальный момент повреждения}
    \label{pict_nvchz:REV_S_a}
\end{subfigure}
\hfill
\begin{subfigure}{0.47\textwidth}
    \centering
    \includegraphics[width=\textwidth]{\fbpath/01.01 ВЧ/561.0771 НВЧЗ/img/REV_S_b.pdf}
	\caption{б) каскадное отключение параллельной линии}
    \label{pict_nvchz:REV_S_b}
\end{subfigure}
\caption{Работа НВЧЗ в режиме реверса мощности}
\end{figure}

\medskip

\item
\textls[-20]{Однако при каскадном отключении линии W2 (выключатель Q3 отключится первым) происходит изменение направления тока и возникает реверс мощности.
Во время переходного процесса может возникнуть ситуация, при которой отключающие ИО полукомплекта А1 срабатывают, а полукомплекта А2 ещё не сброшены (см. рисунок \ref{pict_nvchz:REV_S_b}).
Отсутствие ВЧ-сигнала в канале связи приведёт к отключению неповрежденной линии.
Для исключения такого режима сигнал пуска ВЧ-передатчика от блокирующих органов продлевается на время, задаваемое уставкой \textbf{«Тпрод блок. НВЧЗ»}, если пуск был в течение времени, определяемого уставкой \textbf{«Тср. блок. НВЧЗ»}.}

\item
Пуск ВЧ-передатчика от блокирующих органов снимается при срабатывании отключающего РС или отключающего ОНМОП.

\item
В устройстве предусматривается возможность ручного пуска ВЧ-передатчика через входной сигнал ФСУ \textit{«Ручной пуск ВЧ»}.
Ручной пуск разрешается только при отсутствии срабатывания блокирующих органов (кроме ИО U2 блок НВЧЗ) или отключающего реле сопротивления.

\item
Кроме того, пуск ВЧ-передатчика осуществляется в следующих случаях:
\begin{itemize} 
\item при любых операциях с выключателями (при оперативных переключениях и в циклах АПВ) от входных сигналов ФСУ \textit{«Коммутация В1»} и \textit{«Коммутация В2»} с выдержкой времени на возврат, равной 250 мс.
Это необходимо для исключения излишней работы защиты из-за разновременности включения фаз выключателя;
\item при срабатывании блокировки при сквозных токах через ошиновку, если введен программный ключ \textbf{«БСТО НВЧЗ»};
\item \textls[-10]{при выявлении неисправности в цепях напряжения, если введен программный ключ \textbf{«Пуск ВЧ при БНН»};}
\item от внешней команды пуска ВЧ-передатчика (входной сигнал ФСУ \textit{«Внешний пуск ВЧ»});
\item при обнаружении неисправности канала связи устройством автоматической проверки канала (АПК), если введен программный ключ \textbf{«Блок. НВЧЗ от АПК»};
\item при обнаружении неисправности приемопередатчика, если введен программный ключ \textbf{«Блок. НВЧЗ от Неиспр. ПП»};
\item при обнаружении неисправности терминала системой самодиагностики, если введен программный ключ \textbf{«Пуск ВЧ при неиспр. терм.»};
\item \textls[-35]{при выводе терминала или функции НВЧЗ из работы, если введен программный ключ \textbf{«Пуск ВЧ при выводе»}.}
\end{itemize}

\item
Пуск ВЧ-передатчика снимается при появлении сигнала останова ВЧ-передатчика (логический сигнал \textit{«Запрет пуска ВЧ»}).
Это обеспечивает надежное срабатывание защиты противоположного конца линии при каскадном отключении.
Сигнал \textit{«Запрет пуска ВЧ»} формируется в следующих случаях:
\begin{itemize}
\item при срабатывании внутренних защит и УРОВ (входной сигнал ФСУ \textit{«Останов ВЧ»});
\item при приеме внешнего сигнала \textit{«Останов ВЧ внеш.»}.
Останов ВЧ-передатчика от данного сигнала блокируется в случае срабатывания ДЗШ (входной сигнал ФСУ \textit{«Прием Ср. ДЗШ»}), действующего на отключение выключателя защищаемой линии.
Блокировка продлевается на время, равное 200 мс. 
\end{itemize}

\item
Однако пуск ВЧ-передатчика при неисправности терминала, неисправности канала связи или приемопередатчика, а также при выводе терминала или функции НВЧЗ из работы не блокируется от сигнала \textit{«Запрет пуска ВЧ»}. 

\item
Функциональная схема управления ВЧПП (НВЧЗ) представлена на рисунке \ref{pict_nvchz:VChPP}.

\medskip

\begin{figure}[H]
\centering
\includegraphics[width=0.97\textwidth,height=0.97\textheight,keepaspectratio]{\fbpath/01.01 ВЧ/561.0771 НВЧЗ/img/VChPP.pdf}
\captionof{figure}{Функциональная схема управления ВЧПП (НВЧЗ)}
\label{pict_nvchz:VChPP}
\end{figure}

\medskip

\item
Устройство АПК, входящее в состав ВЧПП, используется для автоматической проверки ВЧ-канала связи.
К устройству от ВЧПП подводятся сигналы \textit{«Контакт АПК»} и \textit{«Неиспр. ПП»}.
Наличие сигнала \textit{«Контакт АПК»}, свидетельствует об отсутствии неисправности в канале связи.
При выявлении неисправности канала связи формируется логический сигнал \textit{«Сигн. Неиспр. КС»}.
При ручном пуске ВЧ-передатчика, срабатывании блокирующих органов или отключающего реле сопротивления формируется сигнал \textit{«Запрет АПК»}, блокирующий действие АПК.
Действие АПК также может быть выведено оперативно от внешнего сигнала \textit{«Вывод АПК»}.
При обнаружении неисправности приемопередатчика формируется логический сигнал \textit{«Сигн. Неиспр. ПП»}.
При введенных положениях программных ключей \textbf{«Блок. НВЧЗ от АПК»} и \textbf{«Блок. НВЧЗ от Неиспр. ПП»} дополнительно формируется сигнал \textit{«Блок. от ВЧПП»}, который действует на вывод НВЧЗ из работы и пуск ВЧ-передатчика.

\item
Функциональная схема контроля ВЧ-канала связи представлена на рисунке \ref{pict_nvchz:Ctrl VCh NVChZ}.

\medskip

\begin{figure}[H]
\centering
\includegraphics[width=0.95\textwidth,height=0.95\textheight,keepaspectratio]{\fbpath/01.01 ВЧ/561.0771 НВЧЗ/img/Ctrl VCh NVChZ.pdf}
\captionof{figure}{Функциональная схема контроля ВЧ-канала связи}\label{pict_nvchz:Ctrl VCh NVChZ}
\end{figure}

\medskip

\item
Обеспечивается формирование сигнала \textit{«Вызов»} при приеме ВЧ-сигнала, длительность которого превышает 5 с.

\item 
Уставочные параметры функциональной схемы управления ВЧПП (НВЧЗ) и схемы контроля ВЧ-канала связи представлены в таблице \ref{app_set:nvchz}.

\end{enumerate}


%%%%%%%%%%%%%%%%%%%%%%%%%%%%%%%%%%%%%%%%%%%%%%%%%% ОТКЛЮЧАЮЩИЕ ИО %%%%%%%%%%%%%%%%%%%%%%%%%%%%%%%%%%%%%%%%%%%%%%%%%%
\item Отключающие измерительные органы

\begin{enumerate}

\item
При несимметричных КЗ подготовка цепи отключения осуществляется от ОНМОП, работа которого подтверждается срабатыванием органов по току обратной последовательности или по приращениям токов прямой и обратной последовательностей, а также органа по напряжению обратной последовательности.
Отключающие органы по приращениям токов прямой и обратной последовательностей рекомендуется применять в сетях с тяговой нагрузкой, для которых характерны значительные небалансы токов обратной последовательности.
Время возврата этих органов составляет 150 мс.

\item
\textls[-10]{Угол максимальной чувствительности ОНМОП равен 110°.
Порог чувствительности ОНМОП определяется уставками срабатывания блокирующих органов по току и напряжению обратной последовательности.}

\item
При симметричных КЗ подготовка цепи отключения осуществляется от отключающего РС, для которого предусмотрена отдельная логика блокировки при качаниях.

\item
Характеристика срабатывания отключающего РС представлена на рисунке \ref{pict_nvchz:RS_otkl_NVChZ}.

\medskip

\begin{figure}[H]
\centering
\includegraphics[width=0.3\textwidth,height=0.3\textheight,keepaspectratio]{\fbpath/01.01 ВЧ/561.0771 НВЧЗ/img/RS_otkl_NVChZ.pdf}
\captionof{figure}{ Характеристика срабатывания отключающего РС}\label{pict_nvchz:RS_otkl_NVChZ}
\end{figure}

\medskip

\item
Ввод в работу отключающего РС (логический сигнал \textit{«Ввод РС откл от БК»}) осуществляется от срабатывания блокирующих органов по приращениям токов прямой или обратной последовательностей или пускового органа по току обратной последовательности с торможением от тока прямой последовательности «ИО I2т пуск НВЧЗ» на время, задаваемое уставкой \textbf{«Тввода РС откл»}.
Факт формирования разрешающего сигнала \textit{«Ввод РС откл от БК»} подхватывается RS-триггером на время, определяемое уставкой \textbf{«Тбл повт. ввода РС откл»}, в течение которого запрещается повторный ввод отключающего РС.
Блокировка повторного ввода отключающего РС может быть сброшена при отключении выключателя при введенном программном ключе \textbf{«Сброс БК НВЧЗ от РПО»}.

\item
Отключающее РС имеет направленную характеристику срабатывания, которая задается уставками \textbf{«Rср РС откл НВЧЗ»}, \textbf{«Хср РС откл НВЧЗ»} и \textbf{«Ф1 РС откл НВЧЗ»}.
Углы наклона характеристики срабатывания органа направленности, обеспечивающего направленность РС, задаются уставками \textbf{«Ф2 ОН РС откл НВЧЗ»} и \textbf{«Ф3 ОН РС откл НВЧЗ»}.
Для отстройки от нагрузки предусмотрен вырез в характеристике срабатывания, который определяется уставками \textbf{«Rнг РС откл НВЧЗ»} и \textbf{«Фнг РС откл НВЧЗ»}.
Работа отключающего РС может быть заблокирована при выявлении неисправности в цепях напряжения при введенном программном ключе \textbf{«БНН РС откл НВЧЗ»}.

\item
Отключающий орган по току обратной последовательности с торможением от тока прямой «ИО I2т откл НВЧЗ» используется при работе на длинных линиях, питаемых мощной системой, в случае недостаточной чувствительности отключающего органа по напряжению обратной последовательности.
Действие «ИО I2т откл НВЧЗ» на отключение разрешается только при отсутствии срабатывания отключающего органа по напряжению обратной последовательности.

\item
Характеристика срабатывания «ИО I2т пуск НВЧЗ», используемого в логике БК, и «ИО I2т откл НВЧЗ» в общем виде показана на рисунке \ref{pict_nvchz:IO_I2t}.
Порог срабатывания горизонтального участка задается соответствующими уставками \textbf{«Iср I2т пуск НВЧЗ»} и \textbf{«Iср I2т откл НВЧЗ»}.
Наклонный участок характеристики проходит через точку, равной  , под углом к горизонтальной оси, соответствующим коэффициенту торможения.
Коэффициент торможения является общим для обоих органов и задается уставкой \textbf{«Кт I2т НВЧЗ»}.  

\medskip

\begin{figure}[H]
\centering
\includegraphics[width=0.5\textwidth,height=0.5\textheight,keepaspectratio]{\fbpath/01.01 ВЧ/561.0771 НВЧЗ/img/IO_I2t.pdf}
\captionof{figure}{Характеристика срабатывания ИО тока обратной последовательности с торможением от тока прямой последовательности}
\label{pict_nvchz:IO_I2t}
\end{figure}

\medskip

\item
После срабатывания отключающих органов осуществляется действие НВЧЗ на отключение выключателя.
Для согласования разновременности срабатывания блокирующих органов полукомплектов защиты на противоположных концах линии предусмотрена выдержка времени \textbf{«Тсогл НВЧЗ»} на подготовку цепи отключения после срабатывания отключающих органов.
Выдержка времени на срабатывание защиты \textbf{«Тср НВЧЗ»} выполнена с функцией сброса набора времени при приеме блокирующего ВЧ-сигнала и останова набора времени при кратковременных помехах в ВЧ-сигнале.
Допустимая длительность указанных помех составляет 2 мс.

\item
Для предотвращения ложного срабатывания НВЧЗ при КЗ за трансформатором ответвления предусмотрен функциональный блок отстройки от повреждений за трансформаторами ответвлений (ОПТО). 
При введенном программном ключе \textbf{«Контр. НВЧЗ от ОПТО»} подготовка цепи отключения разрешается только при наличии подтверждения о КЗ на защищаемой линии от ОПТО.

\item
Для обеспечения надежного действия защиты при включении линии на короткое замыкание или при неуспешном АПВ в устройстве предусмотрена функция ускорения при включении выключателя.
В этом режиме защита действует на отключение без учета наличия блокирующего ВЧ-сигнала в канале связи и без выдержки времени.
Так как при включении выключателя осуществляется пуск ВЧ-передатчика (от входных сигналов ФСУ \textit{«Коммутация В1»} и \textit{«Коммутация В2»}, что описано выше), действие с ускорением возможно только с включаемого конца защищаемой линии.
Ускорение защиты вводится программным ключом \textbf{«Уск. НВЧЗ при вкл. В»}.
Действие на отключение осуществляется при срабатывании блокирующего или отключающего РС либо отключающих токовых органов при наличии сигнала разрешения автоматического ускорения \textit{«Ввод АУ»}.
Также дополнительно контролируется подтверждение наличия КЗ на ЛЭП от ОПТО. 

\item
Действие НВЧЗ на отключение подхватывается от срабатывания блокирующего или отключающего РС либо блокирующих токовых органов. 

\item
Действие НВЧЗ на отключение блокируется при следующих условиях:
\begin{itemize} 
\item при срабатывании блокировки при сквозных токах через ошиновку, которая вводится программным ключом \textbf{«БСТО НВЧЗ»};
\item при выводе терминала или функции НВЧЗ из работы;
\item при обнаружении неисправности канала связи устройством автоматической проверки канала (АПК), если введен программный ключ \textbf{«Блок. НВЧЗ от АПК»};
\item при обнаружении неисправности приемопередатчика, если введен программный ключ \textbf{«Блок. НВЧЗ от Неиспр. ПП»};
\item при переводе действия НВЧЗ только на сигнал (от внешнего сигнала \textit{«Перевод на сигнал»}).
\end{itemize}

\item
Функциональная схема НВЧЗ представлена на рисунке \ref{pict_nvchz:NVChZ}.

\medskip

\begin{figure}[H]
\centering
\includegraphics[width=0.95\textwidth,height=0.95\textheight,keepaspectratio]{\fbpath/01.01 ВЧ/561.0771 НВЧЗ/img/NVChZ.pdf}
\captionof{figure}{Функциональная схема НВЧЗ}
\label{pict_nvchz:NVChZ}
\end{figure}

\medskip

\item 
Уставочные параметры функциональной схемы НВЧЗ представлены в таблице \ref{app_set:nvchz}.

\end{enumerate}

\end{enumerate}